Aide au calcul pour cet exercice : $8^{3}= 512 ~;~ 12 \times 8^{2}= 768$

\medskip

Dans une ville, l'évolution d'une épidémie en fonction du temps écoulé $x$, exprimé en semaines depuis le début de l'épidémie, est modélisée par une fonction $f$ définie sur l'intervalle $[0~;~12]$.

Pour tout nombre réel $x$ compris entre 0 et $12,\: f(x)$ désigne le nombre de malades observés à l'instant $x$.

La courbe représentative de $f$, notée $\mathcal{C}_f$, est donnée dans le repère du plan ci-dessous.

\begin{center}
\psset{xunit=0.85cm,yunit=0.027cm,arrowsize=2pt 3}
\begin{pspicture*}(-0.9,-40)(13,290)
\psaxes[linewidth=1.25pt,labelFontSize=\scriptstyle,Dy=40]{->}(0,0)(0,0)(13,290)
\multido{\n=0+1}{14}{\psline[linewidth=0.15pt](\n,0)(\n,280)}
\multido{\n=0+40}{8}{\psline[linewidth=0.15pt](0,\n)(13,\n)}
\psplot[plotpoints=2000,linewidth=1.25pt,linecolor=red]{0}{12.024}{x dup mul 12 mul x 3 exp sub 4 add}
\uput[ul](6.5,240){\red $\mathcal{C}_f$}
\uput[d](10,-15){Temps écoulé (en semaines)}\uput[r](0,285){Nombre de malades}
\end{pspicture*}
\end{center}

\begin{enumerate}
\item Lecture graphique
	\begin{enumerate}
		\item Avec la précision permise par le graphique, lire $f(5)$.

		Interpréter ce résultat dans le contexte de l'exercice.

		\item Déterminer graphiquement la semaine où a eu lieu le pic de l'épidémie.

		Avec la précision permise par le graphique, déterminer le nombre de malades cette semaine-là.
	\end{enumerate}

\item On admet que la fonction $f$ représentée ci-dessus est définie par :
\[f(x)=-x^3+12 x^{2}+4\]

	\begin{enumerate}
		\item Calculer $f(1)$.

		Interpréter ce résultat dans le contexte de l'exercice.

		\item Vérifier que pour tout réel $x$ de l'intervalle $[0~;~12]$,
		\[f'(x)= 3 x(- x + 8).\]
		\item Compléter le tableau donné ci-dessous.
		
		
\begin{center}
\begin{tblr}{|l|p{7cm}|}\hline
\centering$x$ & 0  \hfill12 \\ \hline
Signe de $3 x$ &  \\\hline
Signe de $(-x + 8)$ &  \\\hline
Signe de $f'(x)$ &  \\ \hline
~&~\\
Variations de $f$ &  \\
~&~\\ \hline
\end{tblr}
\end{center}
		\item Déterminer la valeur exacte du nombre de malades la semaine où a eu lieu le pic de l'épidémie.
	\end{enumerate}
\end{enumerate}
